% Generated draft: requires mathematical and source review.
\documentclass[11pt]{article}
\usepackage[margin=25mm]{geometry}
\usepackage{fontspec}
\setmainfont{Latin Modern Roman}
\usepackage{amsmath,amssymb,mathtools,unicode-math}
\setmathfont{Latin Modern Math}
\usepackage{enumitem,needspace,xurl,hyperref}
\setlength{\parindent}{0pt}
\setlength{\parskip}{5pt}
\setlength{\emergencystretch}{3em}
\newcommand{\N}{\mathbb{N}}
\newcommand{\Z}{\mathbb{Z}}
\newcommand{\Q}{\mathbb{Q}}
\newcommand{\R}{\mathbb{R}}
\newcommand{\C}{\mathbb{C}}
\newcommand{\F}{\mathbb{F}}
\newcommand{\A}{\mathbb{A}}
\newcommand{\PP}{\mathbb{P}}
\newcommand{\E}{\mathbb{E}}
\newcommand{\eps}{\varepsilon}
\newcommand{\dd}{\,\mathrm{d}}
\begin{document}
\section*{UnsolvedMath ID: 701284}
\textbf{Source problem code:} catalogue entry
\subsection{Definitions and mathematical statement}
Let \(n\in\Z\) with \(n\geq 2\), and let \(F_n\) be a free group of rank \(n\). Write \(A_n=\operatorname{Aut}(F_n)\). For \(g\in F_n\), define the inner automorphism \(c_g\) by \(c_g(x)=gxg^{-1}\) for every \(x\in F_n\), and put
\[
I_n=\operatorname{Inn}(F_n)=\{c_g:g\in F_n\}\leq A_n.
\]
Work in the first-order language of groups \(\{\cdot,{}^{-1},1\}\). A subset \(D\subseteq A_n\) is definable with parameters if there exist a nonnegative integer \(m\), a tuple \(\bar a\in A_n^m\), and a first-order formula \(\varphi(x,\bar y)\) such that
\[
\forall\alpha\in A_n,\qquad \alpha\in D\ \Longleftrightarrow\ A_n\models\varphi(\alpha,\bar a).
\]
Parameter-free definability requires \(m=0\).

The question attributed to V. Tolstykh is: for every finite rank \(n\geq 2\), is \(I_n\) a first-order definable subgroup of \(A_n\)? The source does not specify whether parameters are permitted, so the parameter-free and parameter-allowed readings remain unresolved qualifications of this draft. No requirement that one formula work uniformly for all ranks is stated.

An element \(p\in F_n\) is primitive if it belongs to some free basis of \(F_n\). The source additionally states that the set
\[
P_n=\{c_{p^k}:p\in F_n\text{ is primitive},\ k\in\Z\}
\]
is definable in \(A_n\). Here powers are interpreted as integer powers, including the zeroth power. This auxiliary definability assertion is supplied background; its parameter convention and a supporting proof or separate citation are not supplied.
\subsection{Short English statement}
For each finite rank \(n\geq 2\), V. Tolstykh asks whether the inner automorphisms can be defined by a first-order formula inside the group \(\operatorname{Aut}(F_n)\). The source leaves the use of parameters unspecified and reports that inner automorphisms induced by powers of primitive elements form a definable set.
\subsection{Sources}
\begin{itemize}
\item Local UnsolvedMath registry, supplied complete source record: permanent ID \(701284\), exact source problem code \texttt{catalogue entry}, title \emph{Kourovka Notebook Problem 16.93}; upstream revision \texttt{d371ee917ac9647fbb06cd6fc529deec6f3dda82}.
\item \emph{The Kourovka Notebook (Unsolved Problems in Group Theory)}, problem 16.93, proposed by V. Tolstykh in the 16th issue (2006). The supplied background also identifies the 21st issue edition (2026), arXiv identifier \texttt{1401.0300v48}. The supplied collection description names E. I. Khukhro and V. D. Mazurov as editors.
\item The registry labels this problem \texttt{open}. This is a supplied status label, not an independently verified assessment. The cited edition and the auxiliary definability assertion have not been independently checked in preparing this statement-only draft.
\end{itemize}
\end{document}
